\SourceLeaf{196}
\Needspace{8\baselineskip}
\pdfbookmark[1]{Erchinger, Siebenzehneck. Anzeige 1825 December 19}{notice-1825-12-19}
\begin{center}
\rule{.86\linewidth}{.3pt}\par
Göttingische gelehrte Anzeigen.\quad 1825 December 19.\par
\rule{.86\linewidth}{.3pt}
\end{center}
Der Königl. Societät ist abseiten des Herrn \textsc{Erchinger} zu Thuningen im Königreich Würtemberg eine kleine Abhandlung vorgelegt worden, welche die
\begin{center}\emph{Geometrische Construction des regelmässigen Siebenzehnecks}\end{center}
zum Gegenstande hat. Die Allgemeine Theorie der regelmässigen Vielecke hat bekanntlich durch die innige Verbindung, in welche sie mit der höhern Arithmetik gebracht ist, eine neue Gestalt und Erweiterung erhalten; ein, wenn gleich verhältnissmässig nur kleiner Theil derselben ist die Theorie derjenigen Vielecke, die sich geometrisch beschreiben lassen. Seit dem Zeitalter der Griechen wusste man, dass das Dreieck, Fünfeck, Funfzehneck und alle diejenigen Vielecke, welche durch Verdopplung oder wiederholte Verdopplung der Seitenzahl aus diesen entspringen, jene Eigenschaft haben, und man glaubte, behauptete auch wohl ausdrücklich, dass dieses die einzigen seien. Die höhere Arithmetik hat gelehrt, dass dieses ein Irrthum war: indem sie die wahren Quellen der ganz allgemeinen Theorie offen legte, ergab sich von selbst, dass es ausser den genannten Vielecken noch unzählige andere gibt, die geometrisch construirt werden können, von denen das Siebenzehneck das einfachste ist. Die Ueberlegenheit der Analyse, welche das Allgemeinste, wie das Besondere mit gleicher Leichtigkeit umfasst, über die Geometrie, die immer beim Besondern stehen bleiben muss, beim Fortschreiten von den einfachern Fällen zu den zusammengesetztern durch stets vergrösserte Verwicklung aufgehalten wird, und jenen den bekannten nächsten Fall schwerlich jemals ohne fremde Hülfe erreicht hätte, zeigt sich dabei im hellsten Lichte. Inzwischen ist es immer wichtig, interessant und wünschenswerth, dass auch die rein geometrischen Behandlungen fortwährend cultivirt werden, und dass die 
\SourceLeaf{197}
Geometrie wenigstens einen Theil der neuen Felder, die die Analyse erobert, sich aneigne. Ref. ist nicht bekannt, dass bisher jemand die Construction des Siebenzehnecks öffentlich behandelt hätte, ausser Herrn \textsc{Pauker} in den Schriften der Kurländischen Gesellschaft und in seiner Geometrie. Verschieden davon und mehr im rein geometrischen Geiste durchgeführt ist die von Hrn \textsc{Erchinger}, welche in Folgendem besteht. (Die dazu gehörige Figur, eine gerade Linie, auf welcher der Folge nach die Punkte $DBGAIFCE$ liegen, kann jeder sich selbst zeichnen.) Eine nach Gefallen angenommene gerade Linie $AB$ verlängere man rückwärts und vorwärts nach $C$ und $D$ so, dass $AC\times BC=AD\times BD=4AB\times AB$ werden; ferner bestimme man die Punkte $E$, $G$ an beiden Seiten der verlängerten Linie $CA$ so, dass $AE\times EC=AG\times CG=AB\times AB$, und den Punkt $F$ auf der Seite $A$ der verlängerten Linie $BA$ so, dass $AF\times DF=AB\times AB$ wird; endlich theile man $AE$ in $I$ so, dass $AI\times EI=AB\times AF$ werde, wo $AI$ der kleinere, und $EI$ der grössere Abschnitt von $AE$ ist. Man mache dann ein Dreieck, in welchem zwei Seiten jede $=AB$, die dritte $=AI$ wird. Beschreibt man um dieses Dreieck einen Kreis, so wird $AI$ die Seite des in den Kreis beschriebenen regelmässigen Siebenzehnecks sein.

Wenn man die Richtigkeit dieser Construction durch die Vergleichung mit der in den \emph{Disquisitiones Arithmeticae} Art. 354 als ein Beispiel aufgestellter Theorie des Siebenzehnecks prüft, so bemerkt man leicht, dass jene nichts anders ist, als die geometrische Uebersetzung derjenigen Gleichungen, auf welche die Anwendung der allgemeinen Theorie führt: in der That sind die Entfernungen der Punkte $C$, $D$, $E$, $F$, $G$, $I$ von $A$ nichts anderes, als die Grössen, die a. a. O. mit $(8\pdot1)$, $(8\pdot3)$, $(4\pdot1)$, $(4\pdot3)$, $(4\pdot9)$, $(2\pdot1)$ bezeichnet sind, wenn man das positive und negative Zeichen durch die Lage ausdrückt, und die Entfernung des Punktes $B$ von $A$ in eben dem Sinn genommen $=-1$ setzt. Allein das eigentlich Verdienstliche der Abhandlung des Hrn. \textsc{Erchinger} besteht nicht sowohl in der Aufstellung der Construction selbst, da die Analyse bereits den einfachsten Weg vorgezeichnet hatte, als in der rein geometrischen Begründung ihrer Richtigkeit, und diese ist mit so musterhafter mühsamer Sorgfalt, alles nicht rein Elementarische zu vermeiden, durchgeführt, dass sie dem Verf. zur Ehre gereicht, und den Wunsch veranlasst, dass sein in der That nicht gemeines mathematisches Talent alle Aufmunterung finden möge.
